\section{Stability and the exact Dobrushin consequence}
\label{sec:stability-dobrushin}

The center-valued boundary used above is a point in a continuous boundary
space.  This section proves that the obstruction has an open basin and then
identifies the precise Dobrushin statement that follows.

\subsection{A positive-measure family of frustrating boundaries}

For a staple tuple $s=(s_1,\ldots,s_m)$, let
\[
  h_{t,s}(U)=\prod_{i=1}^m p_t(Us_i),
  \qquad
  F_s(U)=\sum_{i=1}^m\dist(\one,Us_i)^2,
\]
and normalize $h_{t,s}\Haar$ to a probability $\gamma_{t,s}$.  Label the
target staple by $m$ and put
\[
 s^0=(\one,\ldots,\one,z^{-1}).
\]

\begin{theorem}[Open stability]\label{thm:open-stability}
Fix $0<\rho<b_*$.  There are an open neighborhood
$\mathcal V$ of $s^0$ in $G^m$, constants $a>0$ and $t_0>0$, such that
\begin{equation}\label{eq:uniform-stability}
 \sup_{s\in\mathcal V}\gamma_{t,s}
 \{U:\dist(\one,Us_m)<\rho\}
 \leq \exp(-a/t),
 \qquad 0<t<t_0.
\end{equation}
Consequently there is an open, positive-product-Haar-measure neighborhood
$\mathcal W$ of the center-valued link boundary in
$G^{B_d^{(1)}\setminus\{e\}}$ on which the conditional bad probability tends
to one uniformly.
\end{theorem}

\begin{proof}
Choose $r$ with $\rho<r<b_*$ and set
\begin{equation}\label{eq:stability-gap}
  \Delta=m(b_*-r)^2>0.
\end{equation}
At $s=s^0$, inequality \eqref{eq:good-lower-bound} implies
\[
 \inf_{\dist(z,U)\leq r}F_{s^0}(U)
 \geq \min F_{s^0}+\Delta.
\]
The map $(s,U)\mapsto F_s(U)$ is uniformly continuous on compact $G^m\times
G$.  Also
\[
 \bigl|\dist(\one,Us_m)-\dist(\one,Uz^{-1})\bigr|
 \leq\dist(s_m,z^{-1}).
\]
Shrink an open neighborhood $\mathcal V$ of $s^0$ so that its compact
closure satisfies
$\dist(s_m,z^{-1})<r-\rho$ and
\begin{equation}\label{eq:energy-perturbation}
 \sup_U|F_s(U)-F_{s^0}(U)|<\Delta/8
 \quad(s\in\mathcal V).
\end{equation}
If $\dist(\one,Us_m)<\rho$, then $\dist(z,U)<r$.  Comparing the constrained
energy with a fixed minimizer of $F_{s^0}$ and using
\eqref{eq:energy-perturbation} leaves a gap of at least $3\Delta/4$ between
the target-good set and $\min F_s$.

The parameter-uniform Laplace bounds in Lemma~\ref{lem:varadhan} convert any
strict energy gap into a uniform exponential probability bound.  Decreasing
the exponent to absorb the heat-kernel and integration errors gives
\eqref{eq:uniform-stability} for some $a>0$ and $t_0>0$.

Finally, the finite family of staple maps is continuous.  The inverse image
of $\mathcal V$ is therefore an open neighborhood $\mathcal W$ of the
constructed link assignment.  Every nonempty open set in a compact Lie group
has positive Haar measure; hence $\mathcal W$ has positive product Haar
measure.  The conditional kernels agree with $\gamma_{t,s(\omega)}$, so the
uniform bad-probability conclusion follows.
\end{proof}

The theorem is stronger than merely choosing a version of a conditional
probability at a null boundary point.  It does not, however, estimate how
likely $\mathcal W$ is under the marginal boundary law as $t\downarrow0$.
That probability contains precisely the energetic information discarded by
a worst-boundary supremum.

\subsection{Dobrushin coefficients}

Let
\begin{equation}\label{eq:coherent-law}
 \nu_t(\dd U)=\widehat Z_t^{-1}p_t(U)^m\,\Haar
\end{equation}
be the conditional law at the all-identity boundary.  It concentrates at
$\one$.  The frustrated law $\gamma_t^z$ concentrates on
$\mathcal M_z$, whose points have distance $L/m$ from $\one$.

\begin{proposition}[Asymptotic conditional separation]
\label{prop:tv-separation}
With the total-variation convention in
\eqref{eq:dobrushin-coefficient},
\begin{equation}\label{eq:tv-limit}
  \lim_{t\downarrow0}\|\nu_t-\gamma_t^z\|_{\TV}=1.
\end{equation}
\end{proposition}

\begin{proof}
Choose $0<r<L/(2m)$ and put $A=B(\one,r)$.  The energy for $\nu_t$ is
$m\dist(\one,U)^2$, with unique minimizer $\one$, so
$\nu_t(A)\to1$.  The closure of $A$ is disjoint from $\mathcal M_z$.
Compactness and the large-deviation upper bound therefore give
$\gamma_t^z(A)\to0$.  Hence
\[
 \|\nu_t-\gamma_t^z\|_{\TV}
 \geq \nu_t(A)-\gamma_t^z(A)\longrightarrow1,
\]
and total variation is at most one.
\end{proof}

\begin{theorem}[Failure of the single-link Dobrushin row criterion]
\label{thm:dobrushin-failure}
Let $j_p$ be the private opposite link associated with $p\ni e$ in
Lemma~\ref{lem:private-opposites}.  Then
\begin{equation}\label{eq:individual-influences}
  \lim_{t\downarrow0}c_{e j_p}(t)=1
  \qquad(p\ni e),
\end{equation}
and consequently
\begin{equation}\label{eq:row-liminf}
  \liminf_{t\downarrow0}\sum_{b\neq e}c_{eb}(t)\geq m.
\end{equation}
In particular, the row-sum condition \eqref{eq:dobrushin-row-condition}
fails for all sufficiently small $t$.
\end{theorem}

\begin{proof}
Fix $p\ni e$.  Begin with the all-identity boundary.  Change only $j_p$ to
the central value, or its inverse, that makes the $p$-staple equal to
$z^{-1}$.  The private-link property leaves all other staples at the
identity.  Thus the two boundary configurations in the supremum defining
$c_{e j_p}(t)$ differ only at $j_p$, and their conditional laws are $\nu_t$
and $\gamma_t^z$.  Proposition~\ref{prop:tv-separation} gives
\[
 \liminf_{t\downarrow0}c_{e j_p}(t)\geq1.
\]
Since every coefficient is at most one, \eqref{eq:individual-influences}
follows.  The $m$ private opposite links are distinct, so summing their
contributions yields \eqref{eq:row-liminf}.  Because $m\geq2$, the row sum is
greater than one for all sufficiently small $t$.
\end{proof}

\begin{remark}[Criterion failure is not nonuniqueness]
Theorem~\ref{thm:dobrushin-failure} says that a particular sufficient
criterion, applied to the redundant microscopic link coordinates, cannot
close at weak coupling.  It neither produces two Gibbs measures nor rules out
uniqueness established by a different metric, a block specification, gauge
fixing, or a method that weights adverse boundaries by their cost.
\end{remark}
